\documentclass[a4paper,12pt]{article}
\usepackage[T2A]{fontenc}
\usepackage[utf8]{inputenc}
\usepackage[russian]{babel}
\usepackage{amsmath,amssymb,mathtools}
\usepackage{geometry}
\geometry{left=25mm,right=24mm,top=23mm,bottom=25mm}
\usepackage{helvet}
\renewcommand{\familydefault}{\sfdefault}
\usepackage{microtype}
\usepackage{booktabs,tabularx,array}
\usepackage{parskip}
\usepackage{needspace}
\usepackage{setspace}
\usepackage{xcolor}
\usepackage{tikz}
\usepackage{hyperref}
\usepackage{fancyhdr}
\definecolor{accent}{HTML}{176A70}
\definecolor{ink}{HTML}{202C34}
\definecolor{soft}{HTML}{657781}
\hypersetup{colorlinks=true,linkcolor=accent,pdfauthor={Лёвин Артём Александрович},pdftitle={Ревью домашней работы — 5 ноября 2026}}
\setlength{\parindent}{0pt}
\setlength{\parskip}{7pt}
\setstretch{1.13}
\setlength{\headheight}{14pt}
\pagestyle{fancy}
\fancyhf{}
\fancyfoot[C]{\footnotesize\color{soft}Лёвин Артём Александрович эксклюзивно для Матвея Горбачёва}
\fancyfoot[R]{\footnotesize\color{soft}\thepage}
\renewcommand{\headrulewidth}{0pt}
\renewcommand{\footrulewidth}{0pt}
\newcommand{\sectionline}{\par\vspace{0.5em}\noindent{\color{accent}\rule{\linewidth}{0.4pt}}\par\vspace{0.8em}}
\newcommand{\smalllabel}[1]{\par\Needspace{5\baselineskip}\addvspace{9pt}\noindent{\small\bfseries\color{accent}#1}\par\nobreak\vspace{1pt}}
\newcommand{\errorlabel}{\textcolor{red!75!black}{\textbf{Ошибка:}}}
\newcommand{\keylabel}{\textcolor{blue!60!black}{\textbf{Ключевая идея:}}}
\newcommand{\resultlabel}{\textcolor{teal!65!black}{\textbf{Итог:}}}
\begin{document}
\thispagestyle{empty}
\vspace*{44mm}
\begin{center}
{\small\color{accent}\MakeUppercase{Математика}}\par
\vspace{15mm}
{\Huge\bfseries\color{ink}Ревью домашней работы}\par
\vspace{7mm}
{\large\color{soft}Производная и касательная}\par
\vspace{23mm}
{\normalsize Лёвин Артём Александрович}\par
\vspace{3mm}
{\color{soft}5 ноября 2026 г.}\par
\end{center}
\vfill
\begin{center}
\begin{tikzpicture}
\draw[accent!55,line width=.75pt] (-5,0) .. controls (-2.5,.55) and (-1,-.25) .. (0,.12) .. controls (1.5,.5) and (2.8,-.38) .. (5,.1);
\end{tikzpicture}
\end{center}
\clearpage
{\LARGE\bfseries\color{ink}Сводная таблица}\par
\sectionline
\renewcommand{\arraystretch}{1.55}
{\small
\begin{tabularx}{\linewidth}{@{}p{13mm}p{27mm}X p{31mm}p{34mm}@{}}
\toprule
\textbf{№} & \textbf{Тема} & \textbf{Суть недочёта} & \textbf{Ваш ответ} & \textbf{Полный ответ}\\
\midrule
\hyperref[task:6]{\textbf{6}} & Монотонность & Не выписаны числовые интервалы. & $-3$, $3$; стрелки на схеме. & $\nearrow$: $(-\infty,-3)$, $(3,+\infty)$; $\searrow$: $(-3,3)$.\\[3pt]
\hyperref[task:9]{\textbf{9}} & Касательная & Ошибка подстановки при $x_0=1$. & $x_0=-1$. & $x_0=-1$; $f(1)=-2$.\\
\bottomrule
\end{tabularx}}
\clearpage
\thispagestyle{empty}
\vspace*{0.42\textheight}
\begin{center}{\LARGE\bfseries\color{accent}Разбор ошибок}\end{center}
\clearpage
\phantomsection\label{task:6}
{\LARGE\bfseries\color{ink}Задание № 6}\par
\sectionline
\smalllabel{Текст задания}
Известно, что $f'(x)=x^2-9$. Укажите интервалы возрастания и убывания функции и абсциссы её точек максимума и минимума.
\smalllabel{Ваше решение}
Вы нашли нули производной: $x=-3$ и $x=3$. Затем верно указали на схеме знаки $+$, $-$, $+$ и нарисовали стрелки возрастания и убывания. Отмечены максимум при $x=-3$ и минимум при $x=3$.
\smalllabel{Ваш ответ}
Абсциссы экстремумов указаны верно. Промежутки монотонности представлены стрелками на схеме, но не выписаны в виде интервалов.
\smalllabel{Ошибка}
\errorlabel\ Неполное оформление ответа: условие прямо просит \emph{указать интервалы}.
\smalllabel{Где именно возникает недочёт}
Последний правильный шаг --- таблица знаков производной:
\[
\begin{array}{c|ccccc}
 x&(-\infty,-3)&-3&(-3,3)&3&(3,+\infty)\\ \hline
 f'(x)&+&0&-&0&+ 
\end{array}
\]
Первый неполный шаг --- завершение решения без явной записи числовых промежутков.
\smalllabel{Почему этого недостаточно}
Стрелки позволяют восстановить ход рассуждения, но отдельная запись интервалов делает ответ однозначным. Числа $-3$ и $3$ делят числовую прямую на три промежутка; на каждом из них знак производной постоянен. Знак производной следует проверять именно \emph{между} этими числами. В точках $-3$ и $3$ производная обращается в нуль, поэтому при записи промежутков, где $f'(x)>0$ или $f'(x)<0$, используют круглые скобки.
\smalllabel{Ключевая идея}
\keylabel\ Если $f'(x)>0$ на промежутке, функция на нём возрастает; если $f'(x)<0$, убывает. Смена знака с плюса на минус соответствует точке максимума, с минуса на плюс --- точке минимума.
\smalllabel{Исправление}
Сохраняем Вашу таблицу знаков и дописываем интервалы возрастания и убывания.
\smalllabel{Полное правильное решение}
\begin{align*}
 f'(x)&=x^2-9=(x-3)(x+3),\\
 f'(x)=0&\iff x=-3\ \text{или}\ x=3.
\end{align*}
При $x<-3$ оба множителя отрицательны, их произведение положительно. При $-3<x<3$ множители разных знаков, их произведение отрицательно. При $x>3$ оба множителя положительны.

Следовательно, функция возрастает на $(-\infty,-3)$ и $(3,+\infty)$; убывает на $(-3,3)$. В точке $x=-3$ знак производной меняется с плюса на минус, поэтому здесь максимум. В точке $x=3$ знак меняется с минуса на плюс, поэтому здесь минимум.
\smalllabel{Проверка}
Подставим по одному значению из каждого промежутка:
\[
 f'(-4)=7>0,\qquad f'(0)=-9<0,\qquad f'(4)=7>0.
\]
Знаки согласуются с таблицей и выводами.
\smalllabel{Итог}
\resultlabel\ Возрастание: $(-\infty,-3)$ и $(3,+\infty)$; убывание: $(-3,3)$; $x_{\max}=-3$, $x_{\min}=3$.
\smalllabel{Задача на закрепление}
Известно, что $g'(x)=x^2-16$. Укажите интервалы возрастания и убывания функции $g(x)$ и абсциссы её точек максимума и минимума.
\clearpage
\phantomsection\label{task:9}
{\LARGE\bfseries\color{ink}Задание № 9}\par
\sectionline
\smalllabel{Текст задания}
Прямая $y=2$ касается графика $f(x)=x^3-3x$. Найдите абсциссу точки касания.
\smalllabel{Ваше решение}
Вы правильно заметили, что касательная горизонтальна, поэтому её угловой коэффициент равен нулю. Получено $f'(x)=3x^2-3$ и найдены кандидаты $x_0=1$, $x_0=-1$. Затем Вы проверили значения функции в этих точках и выбрали $x_0=-1$.
\smalllabel{Ваш ответ}
$x_0=-1$.
\smalllabel{Ошибка}
\errorlabel\ В проверке кандидата $x_0=1$ неверно выполнена подстановка. При вычислении $f(1)$ в рукописной строке используется $-1$ вместо $1$ и получается неверное значение $2$.
\smalllabel{Где именно возникает ошибка}
Последний правильный шаг:
\[
 3x_0^2-3=0\quad\Longrightarrow\quad x_0^2=1
 \quad\Longrightarrow\quad x_0=\pm1.
\]
Первый неверный шаг --- проверка $x_0=1$: в выражении $x^3-3x$ обе буквы $x$ должны быть заменены одним и тем же числом $1$. Фактически при этой проверке смешаны значения $1$ и $-1$.
\smalllabel{Почему это неверно}
После подстановки числа в функцию $f(x)=x^3-3x$ значение аргумента должно быть одинаковым во всех его вхождениях. Например, для $x=1$ первое слагаемое равно $1^3=1$, а второе равно $-3\cdot1=-3$. Значит, $f(1)=-2$, а не $2$.

Важен и логический момент: равенство $f'(x_0)=0$ показывает только горизонтальность касательной. Оно ещё не гарантирует, что эта касательная имеет нужное уравнение $y=2$. Нужно дополнительно проверить $f(x_0)=2$.
\smalllabel{Ключевая идея}
\keylabel\ Для касания графика $y=f(x)$ и прямой $y=c$ в точке с абсциссой $x_0$ одновременно нужны два условия:
\[
 f'(x_0)=0,\qquad f(x_0)=c.
\]
\smalllabel{Исправление}
Оставляем правильно найденные кандидаты $x_0=\pm1$ и повторяем проверку значений функции, подставляя каждый кандидат во все члены формулы.
\smalllabel{Полное правильное решение}
Поскольку у прямой $y=2$ угловой коэффициент равен нулю, в точке касания должно выполняться $f'(x_0)=0$. Поэтому
\[
 f'(x)=3x^2-3,\qquad 3x_0^2-3=0,
\]
\[
 x_0^2=1,\qquad x_0=1\quad\text{или}\quad x_0=-1.
\]
Проверяем первый кандидат:
\[
 f(1)=1^3-3\cdot1=1-3=-2\ne2.
\]
Проверяем второй кандидат:
\[
 f(-1)=(-1)^3-3\cdot(-1)=-1+3=2.
\]
Подходит только $x_0=-1$.
\smalllabel{Проверка}
В найденной точке $f(-1)=2$ и $f'(-1)=3\cdot1-3=0$. Значит, касательная проходит на высоте $2$ и горизонтальна, то есть действительно имеет уравнение $y=2$.
\smalllabel{Итог}
\resultlabel\ $\boxed{x_0=-1}$. Итоговый ответ в тетради правильный; следует исправить проверку альтернативного кандидата $x_0=1$.
\smalllabel{Задача на закрепление}
Прямая $y=16$ касается графика $g(x)=x^3-12x$. Найдите абсциссу точки касания.
\end{document}